\cotefile{13}
\scanpage{29}
\inserted{\deleted{Moyennant conjecture de Hodge,}} On \unsure{récupère} $\unsure{M|\mathcal M}$ comme $\mathrm{Rep}(G)$\inserted{(cf \unsure{a)})}, $T_B$ \inserted{\unsure{(sans structure de Hodge)|(sa structure de Hodge)|(son foncteur de Hodge)}} comme le foncteur oubli sur $\mathrm{Rep}(G)$, $T_{\mathrm{DR}}$ comme $\deleted{\unsure{T_B|T_D}}P\overset{G}{\deleted{\unsure{\otimes|\circ}}\times}(\unsure{T_D|T_B}\otimes_{\mathbb Q}k)$ (cf e), f)), \deleted{la structure} de sorte que $G'\simeq\underline{\mathrm{Aut}}_{\otimes}(T_{\mathrm{DR}})$.

\deleted{Alors $Q$ est un bitorseur \unsure{sous|sur} $(G',G'')$ (cf b)c)d))}

Alors on a $H'\subset G'$, un bitorseur $Q$ sur $(H',H'')$, \deleted{deux \unsure{hom.} $i_1,i_2$ \unsure{d|et} \unsure{avec}} et $T_{\mathrm{Hdg}}\simeq T_{\mathrm{DR}}\overset{H'}{\times}Q$, enfin \deleted{$i_1,i_2\colon G_m\unsure{\mathrel{-}|\longrightarrow}$} $G''\simeq\underline{\mathrm{Aut}}_{\otimes}(T_{\mathrm{Hdg}})$ et $\unsure{i_1|j_1},i_2\colon \unsure{\mathbb G_m|G_m}{}_{k}\unsure{\mathrel{-}|\longrightarrow}G''$ définissant une bigraduation sur $T_{\mathrm{Hdg}}$. Les relations \deleted{\unsure{propr.|prop.}} définissant $H''$ en termes de $i_1$ impliquent que cette bigraduation correspond à une \deleted{\unsure{graduation}} \uline{filtration} \unsure{de|sur} $T_{\mathrm{DR}}$. On en conclut une filtration sur $T_{\mathrm{DR}}\otimes_k\mathbb C$. Or $\varphi$ définit \unsure{un|une} \unsure{isom.} de ce dernier avec $T_B\otimes_{\mathbb Q}\mathbb C$, qui se trouve donc muni d'une filtration.

\uline{Conditions} à imposer aux données a) b) c) d) e) f) g)

\deleted{\illegible{1}} 1) Pour tte repr. (ou une repr. fidèle) $V$ de $G$, la filtration qu'on obtient ainsi sur $V_{\mathbb C}$ \deleted{\unsure{satisfait}} fait de $V$ un module de Hodge\inserted{\unsure{(condition ouverte)}}. 2) Ce dernier est \uuline{motivique}

3) La restriction des scalaires de $\mathbb C$ à $k$ de $V_{\mathbb C}$, déduite de $\varphi\colon T_{\mathrm{DR}}\otimes_k\mathbb C\simeq T_D\otimes_{\mathbb Q}\mathbb C$, correspond à une restriction des scalaires du motif sur $\mathbb C$ défini par la structure de Hodge, grâce à 1). \uline{NB} On ne sait pas expliciter \unsure{2)|3)} et 3) \deleted{\unsure{en termes}} directement en termes des données a) .... !

On trouve alors, \inserted{\uline{si on admet \unsure{la conjecture|les conjectures} de Hodge,}} \unsure{un hom|une tour}
% paper watermark, sideways
{\centering \textsc{Extra Strong}\\
\textit{Ma\unsure{in|nn}\unsure{ol|on}\unsure{uy|ny|y} 1808}\par}
% archive sheet number: 28
